\begin{lemma}\label{lem:smooth-conformal-current}
Let $X$ have complex dimension $n\geq2$, let $\theta$ be a smooth real $(1,1)$-form and let $\gamma$ be Hermitian. Suppose $u,a\in L^1_{\mathrm{loc}}(X)$ satisfy
\[
 \theta+\ddc u=a\gamma
\]
as currents. Then $a$ and $u$ are smooth.
\end{lemma}
\begin{proof}
Taking the exterior derivative and multiplying by $\gamma^{n-2}$ gives
\[
 \mathrm{d}a\wedge\gamma^{n-1}
 =\mathrm{d}\theta\wedge\gamma^{n-2}
      -a\,\mathrm{d}\gamma\wedge\gamma^{n-2}.
\]
At each point, $\xi\mapsto\xi\wedge\gamma^{n-1}$ is an isomorphism from real one-forms to real $(2n-1)$-forms, as one sees in a unitary frame. Its inverse depends smoothly on the point. Thus
\[
 \mathrm{d}a=B-aC
\]
in distributions, for smooth real one-forms $B,C$. Since $a\in L^1_{\mathrm{loc}}$, this first gives $a\in W^{1,1}_{\mathrm{loc}}$. Differentiating the equation and using smoothness of its coefficients gives $a\in W^{k,1}_{\mathrm{loc}}$ for every integer $k$. Local Sobolev embedding makes $a$ smooth. Now $\ddc u=a\gamma-\theta$ is smooth; taking its trace in local coordinates and applying elliptic regularity makes $u$ smooth as well.
\end{proof}

\begin{corollary}\label{cor:smooth-dual-exact-equation}
Let $\omega$ be Hermitian on a connected compact complex manifold of dimension $n\geq2$. Let $0\leq f\in L^1(\mathrm{d}V)$ be nonzero and $b=b_\omega(f)>0$. If $D_{f^{1/n}}$ has a smooth strictly positive maximizer $T\in\mathscr G_\omega$, then $f^{1/n}$ agrees almost everywhere with a smooth nonnegative function. Every threshold subsolution is smooth and satisfies
\[
 \mathrm{MA}_\omega(u)=bf\,\mathrm{d}V.
\]
Such a solution exists, and all threshold subsolutions differ by constants. No root-tail or mass-property assumption is required.
\end{corollary}
\begin{proof}
The attained-threshold theorem supplies a genuine threshold subsolution $u$. The following argument applies to every such $u$. Write $T=\gamma^{n-1}$ for its smooth positive Gauduchon root and set
\[
 a=\left(\frac{bf\,\mathrm{d}V}{\gamma^n}\right)^{1/n}.
\]
The normalization of $T$ is $\int_X\omega\wedge T=1$. Since it maximizes, \cref{prop:threshold-trace-saturation} gives the equality
\[
 A\wedge\gamma^{n-1}=a\gamma^n,
 \qquad A=\omega+\ddc u\geq0.
\]
The positive coefficient measures of $A$ are controlled by this trace, which is absolutely continuous. Thus $A$ has an almost-everywhere positive semidefinite matrix density, with coefficients in $L^n$ because $a\in L^n$.

For every smooth Hermitian form $\beta$, \cref{lem:unbounded-gauduchon-trace-density} also gives
\[
 A\wedge\beta^{n-1}
 \geq\left(\frac{bf\,\mathrm{d}V}{\beta^n}\right)^{1/n}\beta^n.
\]
These pointwise matrix inequalities hold simultaneously almost everywhere for all positive definite matrices: on each coordinate chart first use a countable dense family of constant positive matrices, extended to global smooth Hermitian forms on smaller charts, and then use continuity in the matrix. The matrix arithmetic--geometric mean duality therefore yields
\[
 \left(\frac{A^n}{\gamma^n}\right)^{1/n}\geq a
 \quad\text{almost everywhere}.
\]
Here $A^n$ means only the determinant of its measurable coefficient matrix; no equality with a nonlinear current product has yet been used. The trace equality gives the reverse inequality by arithmetic--geometric mean. Equality forces all eigenvalues relative to $\gamma$ to equal $a$, including where $a=0$. Consequently $A=a\gamma$ as currents.

Apply \cref{lem:smooth-conformal-current} with $\theta=\omega$. Both $a$ and $u$ are smooth, and $f^{1/n}=b^{-1/n}a(\gamma^n/\mathrm{d}V)^{1/n}$ has a smooth nonnegative representative. Only now take the smooth top power: $\omega_u^n=a^n\gamma^n=bf\,\mathrm{d}V$. For two threshold subsolutions, the identity $\omega+\ddc u=a\gamma$ has the same right side. Their smooth difference is therefore pluriharmonic and hence constant on connected compact $X$.
\end{proof}

Smooth dual attainment is thus restrictive even before asking for estimates on the potential. The dimension assumption is used when the exterior derivative of $a\gamma$ determines $\mathrm{d}a$. The corollary does not assert smooth attainment for an arbitrary density or any regularity from a singular maximizing current.

If $f$ is smooth and strictly positive, \cref{thm:gaud-volume-dual-formula} supplies a smooth strictly positive maximizing Gauduchon form from the smooth equation. Thus the corollary recovers that case, including constant-difference uniqueness among all genuine threshold subsolutions.
